% =============================================================================
% Section 3 — Data Model & Mutation Algebra
% =============================================================================
\section{Data Model and Mutation Algebra}
\label{sec:datamodel}

\subsection{The Mutation}

The fundamental unit of state change in \pebble{} is the \emph{mutation}---an
immutable, signed record of a single key-value operation.

\begin{definition}[Mutation]
\label{def:mutation}
A \emph{mutation} is a 7-tuple
\[
  m = \langle \mnid,\; \mseq,\; \mts,\; \mop,\; \mkey,\; \mval,\; \msig \rangle
\]
where:
\begin{itemize}[leftmargin=1.5em]
  \item $\mnid \in \NodeSet$ is the identifier of the node that created $m$.
  \item $\mseq \in \mathbb{Z}^{+}$ is a monotonically increasing counter local
        to $\mnid$: no two mutations from the same node share the same sequence
        number, and sequence numbers only increase across node restarts.
  \item $\mts \in \mathbb{Z}^{+}$ is the wall-clock timestamp (milliseconds
        since the Unix epoch) at the time of creation.
  \item $\mop \in \{\texttt{put},\; \texttt{delete}\}$ is the operation type.
  \item $\mkey \in \KeySet$ is the key being modified.
  \item $\mval \in \Sigma^* \cup \{\bot\}$ is the value: a non-empty UTF-8
        string for \texttt{put}, and $\bot$ (null) for \texttt{delete}.
  \item $\msig \in \{0,1\}^{512}$ is the Ed25519 signature of the node over
        all other fields (see \cref{sec:crypto}).
\end{itemize}
\end{definition}

\begin{definition}[Mutation ID]
\label{def:mutid}
The \emph{mutation ID} of $m$ is the pair $(\mnid, \mseq)$, written
$\mathrm{id}(m) = \mnid\texttt{:}\mseq$.
Mutation IDs are globally unique across the cluster:
\[
  \forall m, m' \in \MutSet:\quad
  \mathrm{id}(m) = \mathrm{id}(m') \implies m = m'.
\]
This is enforced by the sequence collision check in the Store
(see \cref{sec:invariants}, Invariant~\ref{inv:uniqueness}).
\end{definition}

\subsection{Operations and Tombstones}

\pebble{} supports exactly two operations:

\begin{description}[leftmargin=1.5em]
  \item[\texttt{put}$(k, v)$] Associates key $k$ with value $v \in \Sigma^*$.
  \item[\texttt{delete}$(k)$] Creates a \emph{tombstone} mutation for key $k$
        with $\mval = \bot$. The tombstone propagates through gossip like any
        other mutation, preventing deleted data from reappearing on nodes that
        have not yet received the deletion.
\end{description}

\begin{definition}[Live key]
A key $k$ is \emph{live} at node $n$ at time $t$ if the winning mutation for $k$
in $n$'s store has operation \texttt{put}. A key is \emph{deleted} if the winner
is a tombstone. A key is \emph{absent} if no mutation for $k$ has been received.
\end{definition}

\begin{remark}[Tombstone replication]
Tombstones must replicate through gossip exactly like writes. Without tombstone
replication, a node that has received a \texttt{put} but not yet received the
subsequent \texttt{delete} will continue to serve the stale value indefinitely.
\pebble{} stores tombstones in the mutation history and transfers them during
gossip sync. Tombstone GC is a v0.2 concern (\cref{sec:limitations}).
\end{remark}

\subsection{The Mutation History}

Every \pebble{} node maintains an in-memory \emph{mutation history}---a mapping
from mutation IDs to the corresponding full mutation objects:

\[
  H : (\NodeSet \times \mathbb{Z}^{+}) \to \MutSet
\]

The history grows monotonically: once a mutation is added it is never removed
(v0.1). The history enables exact range-based gossip: given a peer's version
vector, the local node can enumerate precisely which mutations the peer is missing.

\subsection{Current State}

The \emph{current state} is the projection of the history onto the winning
mutation per key:

\[
  S : \KeySet \to \MutSet,\quad S(k) = \win\bigl(\{m \in H : m.\mkey = k\}\bigr)
\]

where $\win$ is defined by the LWW total order of \cref{sec:conflict}.
The current state is maintained incrementally: when a new mutation $m$ for key
$k$ arrives, $S(k)$ is updated only if $m$ wins over the existing $S(k)$.

\subsection{Mutation Wire Format}

Mutations are serialised as JSON objects when transmitted over the gossip
protocol or stored in the mutation log. The canonical on-wire representation is:

\begin{lstlisting}[language=json, caption={Mutation JSON representation.}]
{
  "nodeId":    "node-a",
  "sequence":  42,
  "timestamp": 1787060000000,
  "operation": "put",
  "key":       "app:theme",
  "value":     "dark",
  "signature": "3045022100..."
}
\end{lstlisting}

\noindent For \texttt{delete} mutations, \texttt{"value"} is \texttt{null}.
The \texttt{signature} field covers all other fields via the canonical JSON
encoding described in \cref{sec:crypto}.
